\documentclass[10pt,reqno]{amsart}

\usepackage[margin=.8in]{geometry}
\usepackage{amsmath,amssymb,mathtools}
\usepackage{booktabs}
\usepackage{microtype}
\usepackage{tikz}
\usepackage[colorlinks=true,linkcolor=blue,citecolor=blue,urlcolor=blue]{hyperref}
\newcommand{\vTwo}{\nu_2}

\newtheorem{theorem}{Theorem}
\newtheorem{lemma}[theorem]{Lemma}

\newcommand{\N}{\mathbb N}
\newcommand{\Z}{\mathbb Z}
\newcommand{\vTwo}{\nu_2}

\title{A six-dimensional unit-step walk with no collinear triple}
\date{}

\begin{document}

\begin{abstract}
We construct an infinite standard-basis walk in $\N^6$ with no three
collinear vertices.  The proof uses a four-state walk in the Gaussian
integers, two elementary state markers, and a short $2$-adic argument.
\end{abstract}

\maketitle

\section{The construction}

We use $\N=\{0,1,2,\ldots\}$.  Gerver and Ramsey asked whether there is an
infinite sequence of distinct lattice points, with bounded consecutive
distances, no three of which are collinear~\cite{GR}.  The following gives a
direct unit-step construction in six dimensions.

\begin{theorem}\label{thm:main}
There is an infinite sequence $P_0,P_1,\ldots\in\N^6$ such that $P_0=0$,
every difference $P_{n+1}-P_n$ is a standard basis vector, and no three
vertices are collinear.
\end{theorem}

We begin with a four-state Gaussian walk, adapted from the construction in
\cite{CK}.  Define $\sigma_n\in\{0,1,2,3\}$ recursively by $\sigma_0=0$ and,
modulo $4$,
\begin{equation}\label{eq:states}
 \sigma_{2n}\equiv-\sigma_n,
 \qquad
 \sigma_{2n+1}\equiv1-\sigma_n.
\end{equation}
Set
\begin{equation}\label{eq:gaussian}
 u_n=i^{\sigma_n},
 \qquad
 z_n=\sum_{0\le r<n}u_r\in\Z[i].
\end{equation}
Thus $(z_n)$ is a nearest-neighbor walk in the Gaussian lattice.  For
$\varepsilon\in\{0,1\}$, equation~\eqref{eq:states} gives
\begin{equation}\label{eq:dyadic}
 u_{2n+\varepsilon}=i^\varepsilon\overline{u_n},
 \qquad
 z_{2n+\varepsilon}
   =(1+i)\overline{z_n}+\varepsilon\overline{u_n}.
\end{equation}
Indeed, the second identity follows by pairing
$u_{2r}+u_{2r+1}=(1+i)\overline{u_r}$.

For a nonzero rational number $x$, write $\vTwo(x)$ for its $2$-adic
valuation.

\begin{lemma}[Same-state chords]\label{lem:chords}
If $m<n$ and $\sigma_m=\sigma_n$, then $z_m\ne z_n$ and
\begin{equation}\label{eq:chord-valuation}
 \vTwo\bigl(|z_n-z_m|^2\bigr)=\vTwo(n-m).
\end{equation}
\end{lemma}

\begin{proof}
Suppose first that $n-m$ is even.  Then $m=2a+\varepsilon$ and
$n=2b+\varepsilon$ for some $\varepsilon\in\{0,1\}$.  The equality of the
endpoint states and~\eqref{eq:states} imply $\sigma_a=\sigma_b$, and hence
$u_a=u_b$.  The correction terms in~\eqref{eq:dyadic} therefore cancel, so
\begin{equation}\label{eq:halving}
 z_n-z_m=(1+i)\overline{z_b-z_a}.
\end{equation}
Thus halving an even gap preserves equality of the endpoint states and divides
the squared norm of the chord by $2$.

Repeat this reduction $\vTwo(n-m)$ times.  The final gap is odd, so its chord
is a sum of an odd number of Gaussian units.  If that sum is $x+iy$, then each
Gaussian unit has coordinate sum $1$ modulo $2$.  Hence $x+y$ is odd, and
$x^2+y^2$ is odd.  Reversing the reductions proves
\eqref{eq:chord-valuation}; it also shows that the original chord is nonzero.
\end{proof}

\begin{proof}[Proof of Theorem~\ref{thm:main}]
Let $\delta_n\in\{0,1,2,3\}$ represent
$\sigma_{n+1}-\sigma_n\pmod4$.  From~\eqref{eq:states},
\begin{equation}\label{eq:increments}
 \delta_{2n}=1,
 \qquad
 \delta_{2n+1}\equiv-1-\delta_n\pmod4.
\end{equation}
Since $\delta_0=1$, induction gives $\delta_n\in\{1,2\}$ for every $n$.
Consequently, the only possible state transitions are
\begin{equation}\label{eq:transitions}
 01,\ 12,\ 23,\ 30,\ 02,\ 13,\ 20,\ 31.
\end{equation}

We now attach two markers to the four states.  Put
\begin{equation}\label{eq:markers}
 (d_0,d_1,d_2,d_3)=(0,1,i,1+i),
 \qquad
 (c_0,c_1,c_2,c_3)=(0,0,0,1-i),
\end{equation}
and define
\begin{equation}\label{eq:shadow}
 R_n=\bigl(d_{\sigma_n},\ z_n+c_{\sigma_n},\ n\bigr)
 \in\Z[i]^2\times\Z\cong\Z^5.
\end{equation}
If the state changes from $r$ to $s$, then
\begin{equation}\label{eq:shadow-step}
 R_{n+1}-R_n
 =\bigl(d_s-d_r,\ i^r+c_s-c_r,\ 1\bigr).
\end{equation}
The eight transitions in~\eqref{eq:transitions} produce only the six vectors
shown in Figure~\ref{fig:steps}.

\begin{figure}[htb]
\centering
\begin{minipage}[c]{.29\linewidth}
\centering
\begin{tikzpicture}[x=1.25cm,y=1.05cm,every node/.style={font=\small}]
  \draw[gray!70,line width=.6pt] (0,0) rectangle (1,1);
  \fill (0,0) circle (1.5pt) node[below left] {$d_0=0$};
  \fill (1,0) circle (1.5pt) node[below right] {$d_1=1$};
  \fill (0,1) circle (1.5pt) node[above left] {$d_2=i$};
  \fill (1,1) circle (1.5pt) node[above right] {$d_3=1+i$};
\end{tikzpicture}
\end{minipage}\hfill
\begin{minipage}[c]{.67\linewidth}
\centering
\begin{tabular}{c@{\qquad}c}
\toprule
transition(s) & $R_{n+1}-R_n$ \\
\midrule
$01$       & $(1,\ 1,\ 1)$ \\
$12$       & $(-1+i,\ i,\ 1)$ \\
$23$       & $(1,\ -i,\ 1)$ \\
$30$       & $(-1-i,\ -1,\ 1)$ \\
$02,13$    & $(i,\ 1,\ 1)$ \\
$20,31$    & $(-i,\ -1,\ 1)$ \\
\bottomrule
\end{tabular}
\end{minipage}
\caption{The square state marker and the six possible shadow steps.  The first
two coordinates in the table are complex.}
\label{fig:steps}
\end{figure}

We claim that no three points $R_n$ are collinear.  Suppose otherwise, and
choose $a<b<c$ such that $R_a,R_b,R_c$ are collinear.  Write
$p=b-a$ and $q=c-b$.  Since the last coordinate of $R_n$ is $n$, the affine
parameter is fixed:
\begin{equation}\label{eq:slopes}
 \frac{R_b-R_a}{p}
 =\frac{R_c-R_b}{q}
 =\frac{R_c-R_a}{p+q}.
\end{equation}
In the first complex coordinate this says
\begin{equation}\label{eq:convexity}
 d_{\sigma_b}
 =\frac{q}{p+q}d_{\sigma_a}
  +\frac{p}{p+q}d_{\sigma_c}.
\end{equation}
The four numbers $d_0,d_1,d_2,d_3$ are the vertices of a square.  As each is
an extreme point of their convex hull, the strict convex combination in
\eqref{eq:convexity} forces
\begin{equation}\label{eq:equal-states}
 \sigma_a=\sigma_b=\sigma_c.
\end{equation}

The correction terms $c_{\sigma_n}$ now cancel from the second complex
coordinate of every chord in~\eqref{eq:slopes}.  Hence
\begin{equation}\label{eq:gaussian-slopes}
 \frac{z_b-z_a}{p}
 =\frac{z_c-z_b}{q}
 =\frac{z_c-z_a}{p+q}.
\end{equation}
By Lemma~\ref{lem:chords}, for example,
\[
 \vTwo\!\left(\left|\frac{z_b-z_a}{p}\right|^2\right)
 =\vTwo\bigl(|z_b-z_a|^2\bigr)-2\vTwo(p)=-\vTwo(p).
\]
Thus the three nonzero fractions in~\eqref{eq:gaussian-slopes} have
squared-modulus valuations
\begin{equation}\label{eq:slope-valuations}
 -\vTwo(p),\qquad -\vTwo(q),\qquad -\vTwo(p+q),
\end{equation}
respectively.  Equality in~\eqref{eq:gaussian-slopes} would therefore give
\(
 \vTwo(p)=\vTwo(q)=\vTwo(p+q).
\)
This is impossible: if $p$ and $q$ have the same $2$-adic valuation, then
$p+q$ has a larger one.  Thus the shadow $(R_n)$ has no collinear triple.

Finally, list the six shadow steps in Figure~\ref{fig:steps} as
$v_1,\ldots,v_6$.  Starting from $P_0=0$, replace every occurrence of $v_j$
by the standard basis vector $e_j\in\N^6$.  The real-linear map
\[
 T:\mathbb R^6\longrightarrow\mathbb C^2\times\mathbb R,
 \qquad T(e_j)=v_j,
\]
satisfies $T(P_n)=R_n$ for every $n$.  Moreover,
$\sum_{j=1}^6(P_n)_j=n$, so the $P_n$ are distinct.  A collinear triple among
the $P_n$ would map to three distinct collinear points among the $R_n$, which
we have just proved impossible.  This completes the construction.
\end{proof}

\paragraph{Acknowledgment.}
The Gaussian $2$-adic mechanism is adapted from~\cite{CK}.  The two-marker
compression and drafts of this streamlined proof were developed with AI
assistance.  No finite computation is used as a premise.

\begin{thebibliography}{9}

\bibitem{CK}
S.~Cambie and E.~Kalviainen,
\emph{An infinite small-step $\Z^3$-walk with no collinear triple},
arXiv:2609.01766, 2026.

\bibitem{GR}
J.~L. Gerver and L.~T. Ramsey,
\emph{On certain sequences of lattice points},
Pacific J. Math. \textbf{83} (1979), 357--363.

\end{thebibliography}

\end{document}
